\chapter{看懂周期表}

\begin{kaichang}
这份附录不是一章故事，而是一件工具——像卷笔刀旁边的那张乘法口诀表。请把它当作周期表的“使用说明书”：读到正文任何地方，遇到“这个元素在哪一族”“为什么氟比氯更凶”“+2 价是什么意思”这类问题，都可以翻回这里查。

如果你已经读过第 10 章和第 11 章，这里的大部分内容你都见过“故事版”；本附录把它们压缩成可以随手查的表格、坐标和规则。如果你还没读到那里，也没关系——每条规则后面都注明了去哪一章看它的来龙去脉。
\end{kaichang}

\section{先学会读一个格子}

周期表的每一格装着一种元素的“身份证”。不同版本的表排版略有差异，但核心信息只有三样：

\begin{itemize}[itemsep=2pt]
\item \textbf{原子序数}：格子角落的整数，等于原子核里的质子数。它是元素真正的编号——氢 1、碳 6、氧 8、铁 26、金 79。整张表就是按这个数从小到大排的，绝无跳号、绝无重号（第 7 章讲过为什么质子数是身份的底线）。
\item \textbf{元素符号}：格子中央的一到两个字母，第一个大写、第二个小写。$\chem{Co}$ 是钴，$\chem{CO}$ 是一氧化碳——大小写差之毫厘，东西完全不同（符号的读写规则详见附录 A）。
\item \textbf{相对原子质量}：格子下方的长数字，是该元素各种同位素在自然界混合后的平均质量，单位约定省略。氯写作 35.45，不是哪种氯原子的质量，而是两种同位素按比例的平均值。
\end{itemize}

很多表还给格子涂上颜色：金属一种色、非金属一种色、类金属一种色。颜色只是人为的标记，不是元素的真实外貌——金是黄色的，但绝大多数金属其实是银灰色的。

先随手练三次读格子。找 11 号：符号 $\chem{Na}$，钠，食盐里那个“盐味来源”的正离子就是它变的。找 17 号：$\chem{Cl}$，氯，游泳呛水时闻到的消毒水气味的主角。找 26 号：$\chem{Fe}$，铁，你血液里的红色和铁钉生锈的红褐色都跟它有关。三个格子、三种熟悉的生活场景——周期表离你的生活从来不比厨房更远。

\section{坐标系：横行与竖列}

把周期表当地图用，先要会报坐标。

\textbf{横行叫“周期”}，从上往下是第 1 到第 7 周期。同一周期从左到右，电子在同一层“楼层”上逐个增加，性格单调地变化：从急着丢电子的金属，逐渐变成拼命抢电子的非金属。第 1 周期只有 2 格，第 2、3 周期各 8 格，第 4、5 周期各 18 格，第 6、7 周期若把镧系、锕系算回去各有 32 格——这些数字全是电子楼层容量的直接写照（第 9、10 章）。

\textbf{竖列叫“族”}，共 18 列。国际通用的编号就是 1 到 18 族，从左数到右。你还会遇到一套老编号：主族标 A（IA、IIA……VIIIA），副族标 B。两套编号都常见，本书记 1～18 的新编号，见到老编号时对照一下即可。

同族元素最外层电子数相同，所以性格像一个家族：第 1 族（除氢）全是遇水就炸的软金属，第 17 族全是抢电子的好手，第 18 族全是什么都不理的稀有气体。\textbf{查陌生元素的性格，先看它楼上楼下的邻居}——这是周期表最好用的功能，当年门捷列夫就是靠这一招预言未知元素的（第 10 章）。

还有一条从硼（B）斜向下划到砹（At）附近的分界线：线左下方是金属，右上方是非金属，骑在线上的硼、硅、锗、砷、锑、碲等“类金属”两边都沾点。你的手机芯片之所以用硅，靠的正是它这种半金属半非金属的脾气。

\section{四个区块：按“户型”分的街区}

比族更粗的划分是\textbf{区块}——按最后填入的电子住进哪种户型的房间（s、p、d、f 轨道）来分：

\begin{table}[htbp]
\centering
\begin{tabular}{llll}
\toprule
区块 & 位置 & 包含元素 & 共同性格 \\
\midrule
s 区 & 最左 2 列（第 1、2 族） & 氢、碱金属、碱土金属 & 最外层 1～2 个电子，易失电子，性格活泼 \\
p 区 & 最右 6 列（第 13～18 族） & 碳、氮、氧、卤素、稀有气体等 & 金属与非金属都有，生命元素集中于此 \\
d 区 & 中间 10 列（第 3～12 族） & 铁、铜、锌、钛等过渡金属 & 多种价态、有色离子、善当催化剂 \\
f 区 & 表底悬空的镧系、锕系 & 镧系、锕系 & 同排元素性格极其相似，多具放射性（锕系） \\
\bottomrule
\end{tabular}
\end{table}

s 区和 p 区合称\textbf{主族元素}：它们的性格主要由最外层电子数决定，规律最整齐，初中到高中化学打交道最多的就是它们。d 区是\textbf{过渡元素}（过渡金属）：差别藏在次外层，外层几乎一样，所以同一排性格渐变、还常常有好几种价态（第 15 章专门讲）。f 区的镧系、锕系被“流放”到表底，纯粹是为了不让表格宽得印不下（第 10 章）。

一个实用的记法：看一种主族元素在第几族，个位数就是它最外层的电子数。第 14 族的碳、硅最外层 4 个电子，第 16 族的氧、硫最外层 6 个，第 17 族的卤素最外层 7 个。知道最外层电子数，就大致知道了它的价态和脾气。

\section{常见价态速查}

\textbf{化合价}（对离子型化合物也可理解为氧化数）是给成键时的电子“记账”：交出电子记正价，抢来电子记负价。它方便好用，但请记住它只是记账规则，不是原子上真的贴着“+2”的标签。主族元素的常见价态可以从最外层电子数直接读出来：

\begin{table}[htbp]
\centering
\begin{tabular}{llll}
\toprule
族 & 最外层电子数 & 常见价态（离子） & 例子 \\
\midrule
1（除氢） & 1 & $+1$ & $\chem{Na^+}$、$\chem{K^+}$ \\
2 & 2 & $+2$ & $\chem{Mg^{2+}}$、$\chem{Ca^{2+}}$ \\
13 & 3 & $+3$ & $\chem{Al^{3+}}$ \\
14 & 4 & $-4$ 到 $+4$，常共用电子 & $\chem{CH_4}$、$\chem{CO_2}$、$\chem{SiO_2}$ \\
15 & 5 & $-3$、$+3$、$+5$ & $\chem{NH_3}$、$\chem{NO_3^-}$、$\chem{PO_4^{3-}}$ \\
16 & 6 & $-2$、$+4$、$+6$ & $\chem{H_2O}$、$\chem{SO_4^{2-}}$ \\
17 & 7 & $-1$ 为主 & $\chem{Cl^-}$、$\chem{NaCl}$ \\
18 & 8（氦为 2） & 通常 0，极少成键 & $\chem{He}$、$\chem{Ne}$ \\
\bottomrule
\end{tabular}
\end{table}

规律只有一句话：\textbf{主族元素的最高正价一般等于族序数的个位数，最低负价等于“个位数减 8”}。比如第 16 族的氧和硫：最高 $+6$ 价（如硫酸根里的硫），最低 $-2$ 价（如 $\chem{O^{2-}}$、$\chem{S^{2-}}$）。$8$ 这个数字正是“住满一层”的额度——往左丢光外层电子能退回一个满层，往右补齐外层也能凑出一个满层，两条路对应两种价态。但原子不是“为了凑八个”才反应，只是这些排布恰好能量更低（第 10 章的误解警示）。

过渡金属不吃这套口诀：铁有 $+2$、$+3$，铜有 $+1$、$+2$，锰能从 $+2$ 一路变到 $+7$。这是因为它们次外层的 d 电子也能部分参与成键，原因和例子见第 15 章。f 区元素最常见 $+3$ 价，细节超出本书范围。

两个容易绊倒的细节。一是氢：它坐在第 1 族却通常不当金属——遇到非金属时它多半以共用电子的方式记 $+1$ 价（如水、甲烷里），遇到很活泼的金属时反而能抢一个电子记 $-1$ 价（如氢化钠 $\chem{NaH}$）。氢是整张表上最不合群的一格，别拿任何一族的口诀硬套它。二是命名：同种元素有几种正价时，中文用“高、亚”区分——$\chem{Fe^{3+}}$ 叫铁离子，$\chem{Fe^{2+}}$ 叫亚铁离子；补铁药片和铁锈里住的，常常是这两种不同的铁。

\section{趋势一：原子半径——越往右越小，越往下越大}

把周期表当成一张趋势图来读（这里用文字描述这张图）：沿着任意一个周期从左往右走，原子半径\textbf{总体逐渐缩小}；沿着任意一族从上往下走，原子半径\textbf{逐级明显增大}。整张表上，大致是右上角的非金属原子最小，左下角的铯、钫最大。

\textbf{原因}是第 11 章讲的那场拉锯战。同周期内，电子住进同一层楼，层数没变，但原子核里的质子在一个个增加，正电荷越来越强，把同一层电子往怀里拉——于是半径缩水。同族向下，每降一行就新开张一层楼，最外层电子离核远了一大截，还有完整的内层电子像棉被一样把核的吸引力屏蔽掉大半——于是半径膨胀。一句话：\textbf{向右是“同一层楼、核更用力拉”，向下是“楼变高了”}。

先给一组真实的数量级，让“趋势”落地：第二周期从锂到氟，共价半径大约从 \ud{130}{pm} 一路缩到 \ud{60}{pm} 左右——七个元素，半径打了对折还多；而第 1 族从锂到铯，从约 \ud{150}{pm} 涨到 \ud{270}{pm} 上下（金属半径口径）。皮米（pm）是 $10^{-12}$ 米，所有原子的半径都在 $10^{-10}$ 米这个量级：一亿个原子排成一列，大约只有一粒米那么长。趋势讨论的就是这个尺度上百分之几十的差别——听起来小，却足以决定钠遇水爆炸而氖什么都不做。

\textbf{例外与细节}有三处值得知道。第一，d 区和 f 区的半径收缩特别慢：过渡金属一排十个元素，半径几乎没差多少；f 区更是收缩到“镧系收缩”的程度（见本附录的挑战框）。第二，同族从上到下增幅在减小：锂到钠差得很多，铷到铯就差得不多了。第三，稀有气体的“半径”不能直接跟邻居比——见下面的适用范围。

\textbf{适用范围}：原子半径不是像弹珠直径那样绝对固定的数。原子边界本来就是一团概率云，测法不同数值就不同：金属晶体里量的叫金属半径，共价分子里量的叫共价半径，稀薄气体里量的叫范德华半径（通常最大）。稀有气体在多数表里只量得到范德华半径，所以它们的“半径”数字看着比左边邻居还大，这是测量口径不同造成的假象，不算例外。比较半径时，只比同一口径下的数值，并且只取趋势、不抠小数。

\section{趋势二：电离能——越往右越难抢，越往下越好抢}

\textbf{电离能}是从一个气态原子上硬掰下一个电子所需的最低能量，常用的第一电离能就是掰下最外层那个电子的代价。趋势图的方向是：同周期从左往右，第一电离能\textbf{曲折上升}（注意“曲折”二字）；同族从上往下，第一电离能\textbf{逐级下降}。右上角稀有气体最难电离，左下角铯最容易。

\textbf{原因}和半径是同一枚硬币的两面。半径小，最外层电子离核近、感受到的有效吸引力强，自然掰不动——所以半径缩小的方向，正是电离能升高的方向。同族向下，楼高、屏蔽厚，外层电子抓得松，电离能就降。这也是为什么第 1 族从上到下遇水一个比一个猛：钾的外层电子比钠的更好“脱手”（第 12 章）。

\textbf{例外}在短周期里有两处有名的“小台阶”：第 2 族比第 13 族高（铍 > 硼、镁 > 铝），第 15 族比第 16 族高（氮 > 氧、磷 > 硫）。直觉这样理解：硼被掰走的是刚住进 p 房间的那个孤单电子，而铍要拆的是一对已配好对的 s 电子，配对住的更稳；氧那边正好相反——它的 p 房间里有一对电子挤在同一间，彼此排斥，掰走一个反而松快。所以趋势图不是一条光滑斜线，而是带锯齿的上升。

\textbf{适用范围}：电离能是对\textbf{单个气态原子}定义的。它帮你能解释“钠爱丢电子、氯爱抢电子”这类大方向，但不能直接预言某个具体反应快不快、放多少热——那还要看成键放出的能量、溶剂、浓度等一大堆因素（第 27、28 章）。把“电离能小”当成“一遇什么都反应”就过头了。

这条趋势也有它的高光应用。电离能最小的铯，被光一照就能交出电子——老式光电管和光敏传感器就靠这种“见光放电”的脾气工作；反过来，电离能最大的氦和氖怎么轰都很难丢电子，于是霓虹灯里它们只是被电子撞得“兴奋”发光，事后完好如初，一支灯管能亮十几年。同一条趋势，左下角和右上角各有饭碗。

\begin{zhengju}
电子分层住这件事，最好的直接证据就藏在电离能里。对同一个原子连续掰电子：掰第 1 个要一份能量，掰第 2 个要更多（因为剩下的正离子抓得更紧），逐级上升。拿钠（11 个电子）的实验数据看：第一电离能约 \ud{496}{kJ/mol}，第二电离能猛地跳到约 \ud{4560}{kJ/mol}——将近 9 倍！之后第三到第九级平缓爬坡，第十级又是一个陡坎。

这组数字像楼层的剖面图：钠最外面那 1 个电子很好掰；掰掉它之后，下一个电子住在低一层的满员楼层里，代价立刻贵了一个数量级；再往下 8 个电子是同一层，所以只是平缓涨价；最后的 1 个电子住在最贴核的第一层，最贵。其他元素也一样：镁的跳变出现在第 2 级之后，铝在第 3 级之后——\textbf{逐级电离能的“台阶”位置，正好等于该元素最外层的电子数}。20 世纪初的物理学家正是靠这类光谱和电离数据，确认电子不是一锅粥，而是一层一层有定数的。
\end{zhengju}

\section{趋势三：电负性——右上角最“贪”}

\textbf{电负性}衡量的是：原子在\textbf{和别人成键时}，把共用电子对往自己这边拽的本事。它是三条趋势里唯一一条“只在关系中才有意义”的性质。趋势图方向：从左往右电负性增大，从上往下减小——所以整张表的“贪婪之王”是右上角的氟（鲍林标度定为 4.0），越往左下角越“大方”（铯约 0.8）。

\textbf{原因}还是那场拉锯战，只是场景换成了合住。核对电子的吸引力强、原子半径小，共用电子对就被拽得离它更近。氟核的有效吸引力强、半径又小，双料冠军；往左往下，吸引力弱了或距离远了，拽电子的本事就打折。

\textbf{例外}：稀有气体一般不给电负性数值。不是它们“电负性为零”，而是这条性质问的是“成键时谁更能拽”，稀有气体压根不成键，问题本身不成立（氙等少数能被迫成键的元素有特殊标度，那是后话）。

举三个你身边的判断。水分子里氢和氧的电负性差居中，电子对歪向氧——所以水是极性分子，能溶解食盐，却溶不了油（油的碳氢键两边电负性几乎相等，是非极性的），第 21 章会展开这场“相似相溶”的戏。牙膏里的氟化物利用了氟的极端电负性：氟抢电子的劲头比氯还猛，钻进牙齿的羟基磷灰石晶体后抱得更紧，让牙釉质更耐酸。而氯气消毒、碘酒杀菌，靠的也是第 17 族那种逮谁抢谁的强氧化性——同一族的“贪”，剂量不同，既救人也伤人（第 13 章）。

\textbf{适用范围}有三点务必记住。第一，电负性是\textbf{标度值}，不是直接称出来的量：鲍林标度最常见，别的标度数值不同，只比较同一套标度内的相对大小。第二，它用来判断\textbf{共价键里电子对偏向谁}：两原子电负性差得大（如钠和氯），电子几乎整个过户，形成离子键；差得小（如碳和氢），基本平摊；居中（如氢和氧），电子对歪向电负性大的一方，形成极性共价键——水分子里的氢端偏正、氧端偏负，就是这么来的，这也决定了水的一大串怪脾气（第 21 章）。第三，电负性差只是经验判据，键的类型是渐变的连续谱，没有一刀切的边界；更不能拿它去预言反应速率。

\begin{qiaoliang}
三条趋势其实是一条规律的三个侧面，可以压缩成两句话：

\textbf{左右方向}——同一层楼，质子数增加、屏蔽不变，有效吸引力增强：半径减小、电离能增大、电负性增大。

\textbf{上下方向}——楼层增加、屏蔽增厚，外层电子离核远且被遮挡：半径增大、电离能减小、电负性减小。

所以整张周期表存在一个“对角线”式的总趋势：越靠右上角（稀有气体除外），原子越小、抓电子越紧、越贪；越靠左下角，原子越大、电子越松、越大方。金属性与非金属性的分界、各族的活泼性排序，全都是这个总趋势的推论。记趋势不用背六个箭头，记住“两股力的拉锯”加“楼层的涨落”，六个箭头都能现场推出来。
\end{qiaoliang}

\begin{wuqu}
“电负性大的原子，在分子里就带一个完整的负电荷。”——不对。电负性赢的一方只是把共用电子对\textbf{拉得偏近}，电子并没有整个过户。水分子里的氧只是“偏负”（常记作 $\delta^-$），大约相当于零点几个电子的电荷，离完整的 $\chem{O^{2-}}$ 差得远。只有电负性差极大、电子真正整个搬家的组合（如 $\chem{NaCl}$），才谈得上完整的离子电荷。把“偏负”当“负离子”，会错估分子的许多性质——比如水为什么不容易完全电离（第 21、25 章）。
\end{wuqu}

\section{三条趋势合用：三十秒认识一个陌生元素}

这份说明书的真正价值，在于面对一个你从没听过的元素时，能立刻说出个大概。拿锶（\chem{Sr}，第 38 号）练手：

\begin{itemize}[itemsep=2pt]
\item \textbf{定位}：第 5 周期、第 2 族，s 区，楼上楼下是钙和钡。
\item \textbf{猜价态}：第 2 族，最外层 2 个电子，几乎只会是 $+2$ 价，离子 $\chem{Sr^{2+}}$。它的化合物可以直接照抄钙的：氯化锶 $\chem{SrCl_2}$、碳酸锶 $\chem{SrCO_3}$。
\item \textbf{猜脾气}：同族向下电离能减小、金属更活泼——锶应该比钙更急着丢电子，遇水反应更猛，也要像钠一样隔绝空气保存。焰色反应上它给烟花贡献红色，这正是它在生活里最主要的露面方式。
\end{itemize}

再来一个右边的：硒（\chem{Se}，第 34 号）。定位：第 4 周期、第 16 族，氧和硫的楼下。猜价态：最低 $-2$、最高 $+6$，化合物照抄硫的格式。猜脾气：电负性比硫略小、抢电子的劲头稍弱；是人体需要的微量元素，但安全剂量范围很窄——多了中毒，少了生病，这正是第 14 章反复强调“剂量”的原因。你看，一个陌生的名字，三十秒内变成了有坐标、有价态、有脾气的“熟人”。

这就是周期表作为工具的极限与尊严：它给的永远是\textbf{有依据的第一猜想}，不是最终答案。猜想对不对，最后还得看实验——锗被发现前，门捷列夫的“类硅”也只是猜想（第 10 章）。

\begin{shiyan}
\textbf{周期表寻宝游戏}（只需要一张打印的周期表和纸笔，无化学品，安全）

材料：一张带原子序数的周期表（书后或网上打印均可）、纸、笔。

步骤：不看任何说明文字，只凭坐标回答以下问题，然后把周期表翻到有标注的版本核对——
\begin{enumerate}[itemsep=1pt]
\item 找出第 4 周期第 16 族的元素，写出它的符号，并预测它的最低负价。
\item 在钾（\chem{K}）和铷（\chem{Rb}）中，猜哪个遇水反应更剧烈，写出你的理由（一句话）。
\item 找出电负性最大的元素和金属性最强的稳定元素，标出它们的位置，看看是不是一个在最右上、一个在最左下。
\item 随便挑一个你不认识的过渡金属，查它楼下的元素，比较两者的用途（查资料），体会“同族相似”在 d 区是否成立。
\end{enumerate}

观察点：你的预测对了几个？错的那几个，错在“趋势不适用”（比如稀有气体）还是错在推理跳步？把错因写在旁边——这比全部答对更有收获。
\end{shiyan}

\begin{tiaozhan}
\textbf{镧系收缩与对角线相似}（可跳过，不影响工具使用）

第 6 周期的过渡金属（如铪、钽、钨）和它们第 5 周期的同族楼上（锆、铌、钼）半径几乎一样，性格像到难以分离。原因藏在 f 区：镧系 14 个元素把电子塞进很深的 4f 房间，4f 电子对核的屏蔽效果很差，于是镧系一路走下来核的有效吸引力悄悄涨了一截，把后面整个第 6 周期的原子都“勒小”了一圈——这叫镧系收缩。它是趋势图里“同族向下半径增大”这条规律在 d 区几乎失效的原因。

另一个有趣的小规律叫对角线规则：短周期里，某些沿对角线相邻的元素（锂—镁、铍—铝、硼—硅）性格意外相似，因为“向右变小”和“向下变大”两种效应在对角线方向上恰好部分抵消。这两条都不是要背的知识，而是提醒你：周期律是几条竞争因素叠加的结果，到了表格的边角，叠加会出花样——规律依然成立，只是要算得细一些。
\end{tiaozhan}

\section{这份说明书不能替你做什么}

最后交代边界。周期表的趋势和价态表，是\textbf{第一条猜想的生产线}，不是终点：它预测不了反应速率（那是动力学，第 27 章），预测不了反应最终停在哪（那是平衡与自由能，第 28 章），更预测不了分子的三维形状（那是附录 C 和第 19、20 章的事）。遇到具体物质，请把这里的猜想当作提问的起点，顺着五个贯穿问题继续往下挖。

而当你读完元素的故事，想知道这些原子怎样组装成分子、分子又怎样组装成生命时——请翻到附录 C《看懂分子模型》，那是下一件工具。如果某个元素符号、化学式或反应箭头本身让你卡住了，附录 A《化学符号急救箱》负责那一层的救急。两件工具分工明确：本附录回答“这个元素大致是什么脾气、会显出什么价态”，附录 C 回答“它们连起来的分子长什么样、各种画法各保留了什么”。

\section*{练一练}

\begin{sikaoti}
\item 不看表，只凭坐标回答：氯（第 3 周期第 17 族）和溴（第 4 周期第 17 族），谁的原子半径大、谁的第一电离能高、谁的电负性大？每一条各用一句话说理由。
\item 铝的逐级电离能在第 3 级和第 4 级之间出现巨大的跳变。不查数据，画出你预测的铝的逐级电离能“台阶图”（横轴是掰第几个电子，纵轴是能量），并解释跳变的位置说明了什么。
\item 某同学在作业里写：“氧的电负性比氢大，所以水分子里的氧带两个负电荷、两个氢各带一个正电荷。”请指出这句话错在哪一层，并改成正确的说法。
\item 用“三十秒认识陌生元素”的流程处理铷（\chem{Rb}，第 37 号）和碘（\chem{I}，第 53 号）：定位、猜价态、各猜一条物理或化学脾气。然后查资料核对，记录猜错的地方。
\item 为什么“原子半径”这条趋势对稀有气体要打个折扣？如果是你来设计一张更诚实的周期表挂图，你会怎样标注稀有气体的半径，才能不让读者误会？
\item 镓（\chem{Ga}）在铝的楼下、锌的楼上偏右。综合区块、半径和电离能趋势，猜一猜：镓的最高正价是多少？它是金属还是非金属？再查一查它的熔点，你会发现它能在手掌心里熔化——想想这条性质能不能从周期趋势直接推出来，为什么？
\end{sikaoti}
